\documentclass[12pt]{article} % Клас документа, шрифт 12pt

\usepackage[utf8]{inputenc} % UTF-8 кодування
\usepackage[T2A]{fontenc} % Шрифти кирилиці
\usepackage[ukrainian]{babel} % Українська мова
\usepackage{geometry} % Поля сторінки
\geometry{a4paper, margin=2cm} % Формат A4, поля 2 см

\usepackage{xcolor} % Кольори
\usepackage{tcolorbox} % Кольорові рамки
\usepackage{graphicx} % Зображення
\usepackage{amsmath, amsthm, amssymb} % Математика, теореми, символи

\usepackage{caption} % Підписи до рисунків
\captionsetup{labelsep=period} % Роздільник у підписах "Рис. 1."

\newtheorem{theorem}{Теорема} % Середовище для теорем

\usepackage{fancyhdr} % Колонтитули
\pagestyle{fancy} % Використання fancyhdr
\fancyhf{} % Очищення стандартних колонтитулів
\setlength{\headheight}{15pt} % Висота верхнього колонтитулу

\newcommand{\student}{Введіть прізвище та ім'я} % ПІБ студента
\newcommand{\labdate}{\today} % Дата

\fancyhead[L]{\student} % Лівий верхній колонтитул
\fancyhead[R]{\labdate} % Правий верхній колонтитул
\fancyfoot[C]{\thepage} % Нижній колонтитул — номер сторінки

\begin{document}

\begin{center}
    \Large \textbf{Лабораторна робота № 9} \\
    \large \textbf{Тема: Набір математичних текстів та автоматизація в LaTeX}
\end{center}

\section*{Мета роботи}
\textit{навчитися вводити математичні формули та структурувати документ, використовувати кольорові рамки для визначень, теорем та прикладів, а також вставляти зображення}.

\section{Короткі теоретичні відомості}

\begin{tcolorbox}[colback=blue!10!white, colframe=blue!80!black, title=Визначення]
% Введіть визначення квадрата, змініть колір рамки
\end{tcolorbox}

% Введіть текст \\На рисунку ... дивись зразок. 

%додати зображення, файл foto1.png 

\medskip
\textit{Властивості квадрата:}
\begin{enumerate}
    \item Усі кути квадрата прямі.
   %Введіть властивості 2-6 за зразком
\end{enumerate}

%додати зображення, файл foto2.png 

\begin{tcolorbox}[colback=green!5!white, colframe=green!75!black, title=Приклад 1]
% Введіть текст прикладу за зразком
\end{tcolorbox}

%додати зображення, файл foto3.png 

\begin{proof}
\begin{enumerate}
    \item $\angle MAD = \angle ABK = 45^\circ$, а $AD = AB$ як сторони квадрата.
   % Введіть текст доведення 2,3 
\end{enumerate}
\end{proof}

\section{Ознаки квадрата}

\begin{tcolorbox}[colback=blue!5!white, colframe=blue!75!black, title=Ознаки квадрата]
\begin{theorem}
% Введіть текст ознаки квадрата 1) 2)
\end{theorem}
\end{tcolorbox}

\begin{proof}
%Введіть текст доведення використовуючи зразок
\end{proof}

\section*{Контрольні запитання}
\begin{enumerate}
    \item Що таке квадрат? Які його властивості?
    \item Які ознаки квадрата через діагоналі?
    \item Як вставляти зображення та підпис у LaTeX?
    \item Як оформити визначення, теорему та приклад у кольоровій рамці?
\end{enumerate}

\end{document}